\chapter*{Acrónimos}
\addcontentsline{toc}{chapter}{Acrónimos}
Los siguientes acrónimos son empleados a lo largo del documento.
{
\footnotesize
% Configurar el entorno acronym en formato de dos columnas
\begin{acronym}[AAAAAAAA]
\setlength{\itemsep}{0.1pt}
\setlength{\parskip}{0pt}
\setlength{\baselineskip}{10pt}
\setlength{\topsep}{0pt}
\setlength{\partopsep}{0pt}

\acro{AIA}{Atmospheric Imaging Assembly — Conjunto de Imágenes Atmosféricas}
\acro{AR}{Active Region — Región Activa}
\acro{AU}{Astronomical Unit — Unidad Astronómica}
\acro{CACTus}{Computer Aided CME Tracking — Rastreo de CMEs Asistido por Computadora}
\acro{CDAW}{Coordinated Data Analysis Workshop — Taller de Análisis de Datos Coordinados}
\acro{CDS}{Coronal Diagnostic Spectrometer — Espectrómetro de Diagnóstico Coronal}
\acro{CH}{Coronal Hole — Agujero Coronal}
\acro{CME}{Coronal Mass Ejection — Eyección de Masa Coronal}
\acro{DKIST}{Daniel K. Inouye Solar Telescope — Telescopio Solar Daniel K. Inouye}
\acro{DST}{Disturbance Storm Time — Tiempo de Tormenta Perturbadora}
\acro{EFR}{Emerging Flux Region — Región de Flujo Emergente}
\acro{EIS}{EUV Imaging Spectrometer — Espectrómetro de Imágenes en EUV}
\acro{EUV}{Extreme Ultraviolet — Ultravioleta Extremo}
\acro{EUVI}{EUV Imager — Generador de Imágenes en UV Extremo}
\acro{EVE}{EUV Variability Experiment — Experimento de Variabilidad en EUV}
\acro{FIELDS}{Electromagnetic Fields Investigation — Investigación de Campos Electromagnéticos}
\acro{HELCATS}{Heliospheric Cataloguing, Analysis and Techniques Service — Servicio de Catalogación, Análisis y Técnicas Heliosféricas}
\acro{HI}{Heliospheric Imager — Generador de Imágenes Heliosféricas}
\acro{HMI}{Helioseismic and Magnetic Imager — Generador de Imágenes Heliosismológicas y Magnéticas}
\acro{ICME}{Interplanetary Coronal Mass Ejection — Eyección de Masa Coronal Interplanetaria}
\acro{IMaX}{Imaging Magnetograph eXperiment — Experimento de Magnetógrafo de Imágenes}
\acro{IMF}{Interplanetary Magnetic Field — Campo Magnético Interplanetario}
\acro{ISOIS}{Integrated Science Investigation of the Sun — Investigación Científica Integrada del Sol}
\acro{LASCO}{Large Angle and Spectroscopic Coronagraph — Coronógrafo Espectroscópico de Gran Ángulo}
\acro{MC}{Magnetic Cloud — Nube Magnética}
\acro{MDI}{Michelson Doppler Imager — Generador de Imágenes Doppler Michelson}
\acro{MHD}{Magnetohydrodynamics — Magnetohidrodinámica}
\acro{NASA}{National Aeronautics and Space Administration — Administración Nacional de Aeronáutica y del Espacio}
\acro{NFI}{Narrow Field Imager — Generador de Imágenes de Campo Estrecho}
\acro{PFU}{Particle Flux Unit — Unidades de Flujo de Partículas}
\acro{PSP}{Parker Solar Probe — Sonda Solar Parker}
\acro{PUNCH}{Polarimeter to UNify the Corona and Heliosphere — Polarímetro para Unificar la Corona y la Heliosfera}
\acro{SDO}{Solar Dynamics Observatory — Observatorio de Dinámica Solar}
\acro{SECCHI}{Sun Earth Connection Coronal and Heliospheric Investigation	— Investigación Coronal y Heliosférica de la Conexión Sol-Tierra}
\acro{SEP}{Solar Energetic Particle — Partícula Energética Solar}
\acro{SF}{Solar Flares — Llamaradas Solares}
\acro{SILSO}{Sunspot Index and Long-term Solar Observations — Índice de Manchas Solares y Observaciones Solares a Largo Plazo}
\acro{SMEI}{Solar Mass Ejection Imager — Generador de Imágenes de Eyecciones de Masa Solar}
\acro{SOHO}{Solar and Heliospheric Observatory — Observatorio Solar y Heliosférico}
\acro{SOLA}{Subtractive Optimally Localized Averages — Promedios Localizados Óptimamente Sustractivos}
\acro{STEREO}{Solar TErrestrial RElations Observatory — Observatorio de Relaciones Solares-Terrestres}
\acro{Sunspots}{Sunspots — Manchas Solares}
\acro{SW}{Solar Wind — Viento Solar}
\acro{SWEAP}{Solar Wind Electrons Alphas and Protons — Electrones, Alfas y Protones del Viento Solar}
\acro{SWOOPS}{Solar Wind Observations Over the Poles of the Sun — Observaciones del Viento Solar Sobre los Polos del Sol}
\acro{TRACE}{Transition Region and Coronal Explorer — Explorador de la Región de Transición y Coronales}
\acro{UVCS}{Ultraviolet Coronagraph Spectrometer — Espectrómetro Coronógrafo Ultravioleta}
\acro{WDC}{World Data Center — Centro Mundial de Datos}
\acro{WFI}{Wide Field Imager — Generador de Imágenes de Campo Amplio}
\acro{WISPR}{Wide-field Imager for Solar Probe — Generador de Imágenes de Campo Amplio para Solar Probe}

\end{acronym}
}